\documentclass[ecta,nameyear,draft]{econsocart}
\input{preamble}
\externaldocument[][nocite]{build/supp}[supp.pdf]
\externaldocument[][nocite]{build/appendix}[appendix.pdf]
\begin{document}
\begin{frontmatter}
\title{Neural Bellman Operators}
\runtitle{Neural Bellman Operators}
\begin{aug}
\author[id=au1,addressref={add1}]{\fnms{Qian}~\snm{QI}\ead[label=e92ECTA]{qiqian@pku.edu.cn}}
\address[id=add1]{Peking University}
\end{aug}
\begin{abstract}
Neural Bellman Operators construct and learn continuation operators for feasible
decisions under a specified recursive economic objective. We derive explicit
regularity of the original action marginal in a constrained investment economy
with continuous innovations and nested average value at risk, including binding
capacities and moving adverse tails. Paid original-risk value enclosures supply
outward derivative labels. A bounded action-marginal readout has an exactly
checkable primal--dual training gap; feature-aware continuation contrasts and
constrained action bounds state what additional work controls original policy
loss. A recursive level-and-contrast theorem shows why a state-dependent critic
anchor cannot be discarded when exported to an earlier Bellman operator.
Complete proofs retain the own-future construction, stochastic-control,
recursive-utility and boundary assumptions. The prospective independent study
compares two learning losses and acquisition budgets for neural, matched-hinge,
quadratic, polynomial, spline and grid critics against a strong online method,
with all acquisition, refusal, qualification and complete-service costs.
Its source, actor and new state-sample boundaries are fixed separately.
At this source-freeze stage no new numerical outcome is asserted; the completed
author manuscript will report the actually executed original-policy and cost
comparisons, including adverse results and unresolved adoption regions.
\end{abstract}
\begin{keyword}
\kwd{Dynamic programming}\kwd{Neural approximation}\kwd{Recursive risk}
\kwd{Policy qualification}\kwd{Computational economics}
\end{keyword}
\end{frontmatter}
\section{Introduction}
\label{sec:introduction121}

A dynamic decision maker buys an approximation because it changes the cost of
making a good decision. In a constrained investment problem, an accurate
continuation level can still give the wrong incentive to invest, while a
continuation with a large common level error can choose the correct action.
Under recursive risk, the difficulty is greater: changing a continuation can
also change which shocks receive adverse weight. A numerical method must
therefore connect its acquired information, its fitted critic, its feasible
actor and its computational work to the original recursive objective.

Neural Bellman Operators (NBOs) address this chain by constructing and learning
recursive continuation operators together with an account of their policy
error. The object is the Bellman operator of the specified economy. Its
criterion, information filtration, constraints and boundary conditions remain
part of that object throughout approximation and deployment. The paper retains
both a constructive own-future neural operator and an analysis of arbitrary
acquired critics. It develops their links with stochastic control and, under
separate explicit assumptions, nonlinear Brownian recursive utility. A result
for the nested risk investment application is not automatically a welfare
improvement in a different expected-utility or nonlinear recursive-utility
model.

The central computational question is which continuation information is worth
acquiring. Earlier NBO calculations found a concrete mismatch: every
certificate-oriented fit in a retained comparison reduced a residual-span
criterion, but the aggregate number of rejected first actions increased from
644 to 757. The quadratic control remained a strong complete-service baseline.
The subsequent policy-comparison calculations preserved that economic
criterion and constructed safe actual improvements, yet acquisition still
used a value loss. Those results motivate a new learning calculation in the
original model. Their observations and costs retain their historical
identities in the development companion; the present sample does not reuse
them.

We learn the marginal continuation cost after evaluating the original risk
operator. This order matters. In the two-date investment economy, the original
optimal continuation is the last-date Bellman image of the known terminal
function. It is acquired through paid outward calculations. The root critic
retains the known action cost and learns the action-dependent part of the
integrated continuation. A state-only anchor completes its value level.
The resulting actor is proposed by piecewise minimization of that learned
objective over the original capacity interval. Floating critical-point
calculations generate a proposal; admission requires a certificate for the
original complete policy. The same acquired values, features where relevant, feasible
actions and certification rules are available to the conventional controls.

The mathematical contribution has three connected parts. First, we derive a
continuous original-action marginal and explicit moduli from the nonlinear
investment technology, binding capacity and moving AVaR tails. The central
interval excluded from the adverse tail can touch either shock boundary. Its variational inequalities
control movement of the tail without differentiating a numerical sort or an
unconstrained optimizer. Finite original-risk value enclosures then provide
outward derivative labels, with the precision cost of finite differences paid
explicitly.

Second, a bounded readout is fitted to minimize the worst action-marginal discrepancy over
those outward labels. The finite training problem is convex for each declared
feature family. Exact primal and dual bounds certify the actual stored
coefficients without treating a floating solver's success flag as a proof.
An explicit feature-range construction states how to extend the acquired
information to off-design action comparisons. Its integral along the feasible action path
controls policy loss and possible policy changes; a capacity-sensitive
stationarity bound sharpens the account when available. The argument applies
to a nonconvex fitted polynomial or spline through the all-action comparison,
and never assumes that target smoothness removes a fitted ReLU kink. This
extension is a proved construction whose numerical cell bounds and cost must
also be supplied before it can certify an off-design decision. The current
experiment audits the finite-node fits and admits policies through the
separate original-model certificate; it does not claim to have executed that
global feature-cell calculation.

Third, we state what is lost if the critic's anchor is discarded recursively.
In this very investment transition, a state-dependent continuation level that
does not change a last-date action can shift the root marginal by a known
nonzero amount. A full level-and-contrast recursion therefore carries the
anchor error, the action gap and the future policy error separately. A
noncash-invariant recursive preference receives its own continuation
Lipschitz account. It cannot inherit a monetary shift identity merely because
both criteria have a Bellman notation.

Derivative-assisted economic dynamic programming already uses Hermite
approximation with value and slope information \citep{CaiJudd2015}.
Derivative supervision and task-based model learning also have established
precedents \citep{Czarnecki2017Sobolev,Donti2017Task,ElmachtoubGrigas2022}.
The training certificates use convex duality and affine supporting bounds;
these general tools do not themselves supply a claim of methodological novelty.
The development history also contains certificate-aware readouts and an
action-null quotient analysis. The additional results here concern their
explicit construction and verification for the original recursive risk
operator, including tail changes, paid finite-risk information, feasible
actions and recursive value levels. A three-probe lower bound identifies a
limitation of that stated information interface; it is not a lower bound for
every algorithm with access to the investment formulas.

The numerical comparison changes the learning loss before changing the
evaluation sample. Two separately executed acquisition budgets use nested
state designs. Seven stored representations each have a value-loss and an
action-marginal-loss critic; the online original-model method is included at
both risk levels. Quadratic, polynomial, cubic spline, multilinear grid and
semi-analytic controls form the conventional menu. An exactly corresponding
hinge encoding makes the neural representation claim testable: a geometry
shared with that control is not a uniquely neural contribution.

The protocol fixes 64 new independent initial states, complete original-policy
evaluation, five fresh matched timing processes, exact independent
reconstruction, simultaneous positive and adverse gain inference and full
acquisition/service cost comparisons. It also fixes how unresolved and failed
calculations are retained. The empirical question is whether action-marginal
learning changes decisions, rejections and complete work enough to compensate
for its acquisition cost. A reduction of its own finite-node fitting objective
alone does not answer that question.

The main article presents the economic operator, its original-risk
qualification and the new learning construction. The supplement contains
complete proofs and retains the earlier action, tail-information and policy
improvement theorems. The development companion preserves the earlier
applications, numerical records, negative findings and source-bound revision
history. This organization keeps the active economic and mathematical
argument readable while retaining the research needed to assess it.

\input{new121/core_main121}
\input{new/centeredrisk92-active99}
\input{new110/continuous-main}
\input{new119/foundations_main}
\input{new121/marginal_main}
\section{An independent comparison of learning losses}
\label{sec:numerical121}

This section is the prospective manuscript component at source freeze F.
The R121 economic acquisition and independent evaluation have not yet been
executed at that boundary. The completed author source replaces this paragraph
with the actual execution identities, complete numerical tables and economic
interpretation. No earlier experiment supplies R121 outcomes.

\subsection{The original decision and the acquired information}
\label{subsec:design121}

The original model is \eqref{eq:investment95}--\eqref{eq:cost95} at $d=2$,
two dates and risk levels $1/2$ and $3/4$. Investment raises the two next-date
stocks by $a/2$ and $a/4$, whereas the continuous innovation shifts them by
$(z,-z)$. It therefore changes the stocks' imbalance without changing their
aggregate. The decision trades the current adjustment cost against tracking,
the low-stock penalty, capacity and the original risk-sensitive continuation.
All objective quantities below use the original normalized cost units; the
experiment has no monetary or consumption-equivalent calibration.

Two endpoint state grids have
$5^2$ and $9^2$ states. Each buys five normalized-action derivative contexts
and their fixed feasible probes, or thirteen original-value contexts per
state. The protocol requires a separate actual acquisition process for each
budget. Every method
within it receives the same complete information bundle and is charged for
its declared standalone acquisition. The optimized last-date continuation is
acquired from the original known terminal primitives. It is neither a free
oracle nor a target supplied by another representation. The root critic learns
the original integrated continuation difference; its separately stored
state-only anchor completes the value level. The future deployed controller
is unchanged, with $\beta 2^{-20}$ charged once in a complete-policy upper
bound.

\subsection{Representations, losses and the complete decision service}
\label{subsec:methods121}

The seven stored families are a convex quadratic, an anchored ReLU critic,
its exactly corresponding conventional hinge encoding, a tensor cubic
spline, unrestricted polynomials of degrees two and three, and a multilinear
grid. The VALUE and MARGINAL fits use the same coefficient box and features
within a family. VALUE uses every paid value, including the probes. MARGINAL
uses the outward original-action labels. The online baseline computes its
proposal directly from the original model. Every proposal faces the same
complete original-law certificate, refusal, conventional fallback,
integration refinement and rational totalization.

\begin{table}[t]
\centering\small
\caption{Declared learned continuation representations. Counts are the stored
readout coefficients at budgets 5 and 9. Every family receives the same paid
information within a budget; the two losses face the same family and box.}
\label{tab:representations121}
\begin{tabular}{p{0.18\textwidth}p{0.60\textwidth}r}
\toprule
Family & Anchored representation and restriction & Counts\\
\midrule
Quadratic & $v,x_0v,x_1v,v^2$; nonnegative $v^2$ coefficient & $4/4$\\
Polynomial 2 & The same monomials with signed coefficients & $4/4$\\
Polynomial 3 & Total degree at most three; every monomial contains $v$ & $10/10$\\
ReLU & Three base features and 36 fixed anchored hinges; nonnegative nonlinear readouts & $39/39$\\
Hinge & The exact same features and readout, independently encoded in a coordinate loop & $39/39$\\
Cubic spline & Tensor clamped cubic basis; anchored action-axis restriction & $48/100$\\
Linear grid & Tensor linear hats on the actual state and action grids & $100/324$\\
\bottomrule
\end{tabular}
\end{table}

All signed coefficients lie in $[-128,128]$; the stated nonnegative ones lie
in $[0,128]$. The hidden ReLU directions are fixed before acquisition and
are not selected using held-out outcomes. Stored readouts use 30-bit dyadic
rounding. The independently reconstructed finite-box optimization gap applies
to these actual coefficients. It does not assert that a floating optimizer
returned the exact mathematical minimizer. The matched hinge is not an
independent fit; it pays the complete corresponding ReLU phase as well as
its own encoding, identity check and audit.

The learned proposal enumerates the declared action pieces, endpoints and
computed real critical points. Numerical root enumeration proposes an action;
the original economic certificate decides whether to deploy it. A rejected
first proposal remains an outcome of the learning method even when a later
conventional fallback completes the decision. The online control uses a
512-midpoint, 24-iteration original-root proposal. Every completed policy must
meet the original regret target $1/256$ and interval-width target $1/8192$.
Thus a method cannot compensate for an inaccurate actor by silently relaxing
the economic accuracy requirement.

\subsection{Independent reconstruction and the statistical unit}
\label{subsec:validationdesign121}

The source/protocol commit precedes all acquisition. The acquired model
commit precedes a single new operating-system entropy draw. The state commit
precedes every outcome. There are 58 policy cells, 3,712 primary singleton
decisions and 9,280 decisions in five serial fresh timing processes. The
independent economic reconstruction evaluates all distinct final-policy and
raw-proposal keys with the same rational design, and the separate training
audit reconstructs the actual algebraic feature rows and finite-box gaps.
The 64 states, rather than risk cells or repeated clocks, are the independent
statistical units. Their population is the product uniform law on the two
40-bit initial-state coordinates, conditional on the stated independent
OS-entropy assumption. The original innovation distribution remains
continuous. Independent validation reconstructs optimized-continuation values
with exact rational primitive formulas and adds the implemented-future
allowance once; compatibility with a production interval does not by itself
certify that interval's narrower endpoints.

Gain is original cost under VALUE minus original cost under MARGINAL on the
same state. It can have either sign. Both completed policies have regret in
$[0,1/256]$, so a paired gain lies in $[-1/256,1/256]$. The fixed joint
95 percent family in Proposition~\ref{prop:jointinference121} includes the
two optimal-value means, all 58 policy-regret means, all 28 signed loss
contrasts, 112 signed threshold probabilities and both endpoints of all 58
first-rejection probabilities. Identical policy keys cancel exactly without
creating new sample units. Raw proposals have no final-service regret cap;
their complete sample diagnostics are reported separately. Cross-family
sample comparisons and compute-price regions do not acquire an undeclared
population-superiority interpretation.

The complete prospective design and reporting family are
retained in \texttt{PROTOCOL121.md}, \texttt{INDEPENDENT\_PROTOCOL121.md} and
\texttt{REPORTING\_SPEC121.md}.

\section{Economic use and the remaining research task}
\label{sec:discussion121}

The role of a Neural Bellman Operator is to acquire and reuse continuation
information that changes feasible economic decisions at an acceptable total
cost. The marginal construction gives that objective a precise training and
verification account. It also exposes two costs that cannot be removed by
renaming a fitting statistic: integration must track the original risk
operator, and an exported value must retain its state-dependent level.

The acquisition budgets in this experiment change information density and
representation size within the same two-state economy. They do not establish
scalability in horizon, state dimension or innovation dimension. The retained
theorems describe those separate objects and their assumptions; their
historical experiments and unresolved regions remain in the development
companion. The stopping, accuracy-allocation and risk-information questions
therefore remain part of the NBO programme under the original economic
criteria.

A neural-specific advantage requires a distinction that a conventional
method with the same information, features, feasible policies and resources
cannot reproduce. The exactly matched hinge control makes one such boundary
explicit. Conditional adoption additionally requires the original objective
gain to pay for all acquisition and complete-service work at a stated price
and volume. These are empirical and mathematical requirements for an
economically useful NBO, not consequences of a successful optimization status
or a smaller internal loss.

The measured ReLU--hinge comparison needs a further computational distinction.
The two encodings carry identical mathematical features and readouts. Under
the frozen acquisition procedure, the coordinate-loop hinge implementation
pays for the corresponding ReLU fit and its own encoding, identity check and
audit. Those are real costs of the executed implementation, and remain in the
ledger. A conventional hinge representation can also use the same batched
evaluation and elementary operations as the ReLU implementation. Consequently
a cost difference caused by the declared copy or evaluation procedure cannot
establish a neural-specific computational advantage. Any sufficient price
region against the recorded implementations is conditional on that measured
menu. Its interpretation must also retain the other strong conventional
alternatives and the distinction between implementation cost and a limitation
of conventional representation.

\bibliographystyle{ecta-fullname}
\begingroup\emergencystretch=1em
\bibliography{references99,references110-additions,references120,references121}
\endgroup
\end{document}
